\subsection*{§. 190.\\ 双曲空间运动的分类。}

在双曲空间的一个运动中，可能有某个测地曲面在自身上滑动（而不把两面倒置），也可能没有这样的曲面。在第一种情形，空间的运动已由该曲面在自身上的运动完全确定；因此，按照曲面自身运动所属的类别（§. 175），称空间运动为椭圆型、双曲型或抛物型。若没有测地曲面在自身上滑动，则称该运动为螺旋型。

为了从相应线性代换 (15) 的系数 $\alpha,\beta,\gamma,\delta$ 判别运动的类别，先假定代换为幺模的。注意，总可以把代换 (15) 换成它的一个变换，即把给定运动换成一个仿射运动；还要回顾，从代换 (15) 过渡到它的任意变换时，和 $\alpha+\delta$ 是一个不变量，记为 $j$。

现在，若运动中有一个测地曲面在自身上滑动，就以一个仿射运动替换它，使平面 $\eta=0$ 在自身上滑动。相应代换的系数将成为实数；由于复变量 $\xi+i\zeta$ 在平面 $\eta=0$ 上受到同一代换（而 $\xi-i\zeta$ 受到共轭代换），故有
\[
\begin{array}{ll}
\text{椭圆型运动} & \text{当 } |j|=|\alpha+\delta|<2,\\
\text{抛物型运动} & \text{当 } |j|=2,\\
\text{双曲型运动} & \text{当 } |j|>2.
\end{array}
\]
在这三种情形中，代换的不变量 $j$ 都是实数。反过来，若这个不变量是实数，运动便属于这三类之一，即不是螺旋型。事实上，当给定代换的两个不动点互异时，可把它变为形如
\[
 z'=\alpha z \quad(\alpha\delta=1),
\]
的代换；当两个不动点重合时，则变为
\[
 z'=z+a
\]
的代换。在后一情形，Poincaré 公式给出表示欧几里得空间的一个平移；有无穷多个测地曲面（由垂直于极限平面的平行平面表示）在自身上滑动，故运动为抛物型。在前一情形，由于 $\alpha+\delta$ 为实数且 $\alpha\delta=1$，$\alpha,\delta$ 或为实数，或为互相共轭且模 $=1$ 的虚数；第一种情形的运动为双曲型，第二种为椭圆型。

综上：\emph{若不变量 $j=\alpha+\delta$ 为复数，则运动为螺旋型；若 $j$ 为实数，则依次按 $|j|<2$、$|j|=2$ 或 $|j|>2$，运动为椭圆型、抛物型或双曲型。}

再注意，除抛物型情形外，极限平面上总有两个互异的不动点 $A,B$；通过 $A,B$ 且垂直于该平面的圆，表示一条在运动下沿自身滑动的双曲空间测地线（运动的轴）。

于是可把相应代换化为
\[
 z'=\frac{\alpha}{\delta}z,\qquad \alpha\delta=1
\]
由此得到下列情形。
在椭圆型情形，轴上的所有点都保持不动，而空间运动是绕该轴的\emph{旋转}。令
\[
 \alpha=e^{i\varphi},\qquad \delta=e^{-i\varphi},
\]
则
\[
 z'=e^{2i\varphi}z
\]
而旋转的\emph{幅度} \(\Omega\) 由 \(2\varphi\) 给出；又有 \(j=\alpha+\delta=2\cos\varphi\)，因而
\[
\tag{18}
 j=2\cos\frac{\Omega}{2}.
\]
在双曲型情形，运动的轴在自身上滑动，所有包含该轴的测地曲面也随之在自身上滑动；轴上的每一点移动一个常长度 \(d\)（按双曲度量计算），称为双曲运动的\emph{幅度}。这个幅度 \(d\) 仍只依赖不变量 \(j\) 的值。事实上，此时 \(\alpha,\delta\) 为实数且 \(\beta=\gamma=0\)，公式 (17*) 给出
\[
 \zeta'=\frac{\zeta}{\delta^2},
\]
而对于由 \(\zeta\) 轴表示的测地线，有
\[
 ds=\frac{d\zeta}{\zeta},
\]
故
\[
 d=\log\frac{\zeta'}{\zeta}=\pxChThirteenDistanceLog.
\]
\ednote{原书第428页先给出 \(\zeta'=\zeta/\delta^2\)，随后却把 \(\log(\zeta'/\zeta)\) 印作 \(+2\log\delta\)。直接代入前式必得 \(-2\log\delta\)，故本译在正文中采用负号；下一式中的 \(\cosh(d/2)\) 为偶函数，因而不受这个符号修正影响。}
另一方面
\[
 j=\alpha+\delta=\delta+\frac1\delta,
\]
于是得到所求公式
\[
\tag{18*}
 j=2\cosh\frac d2,
\]
它应与 (18) 对照。

最后，在螺旋型情形，运动的轴在自身上滑动，同时空间绕轴旋转；该运动显然由具有同一轴的一个椭圆型运动与一个双曲型运动复合而成。

除了由 (15) 表示的本类运动，还可考察与 (16) 对应的运动，称为\emph{第二类运动}。对这些运动，两个对应图形不再直接全等，而是\emph{反向全等}。关于这些运动的公式显然可由 (17) 中交换 \(z\) 与 \(z_0\) 得到。

第二类运动中特别值得注意的是周期为2的运动，它们称为\emph{反射}，并分成两类。第一类是真正关于一个测地曲面的\emph{对称}；第二类则由这样的反射与绕一条测地轴旋转一个平角复合而成，该测地轴垂直于反射的测地曲面。
